    \documentclass[12pt,a4paper]{article}

    % Encoding and fonts
    \usepackage[utf8]{inputenc}
    \usepackage[T1]{fontenc}
    \usepackage{lmodern}
    \usepackage{textcomp}

    % Page layout
    \usepackage[top=1in, bottom=1in, left=1in, right=1in]{geometry}

    % Typography
    \usepackage{microtype}
    \usepackage{parskip}

    % Language
    \usepackage[english]{babel}

    % Headers and footers
    \usepackage{fancyhdr}
    \pagestyle{fancy}
    \fancyhf{}
    \fancyhead[L]{\leftmark}
    \fancyhead[R]{\thepage}
    \fancyfoot[C]{\thepage}
    \renewcommand{\headrulewidth}{0.4pt}
    \renewcommand{\footrulewidth}{0pt}

    % Section styling
    \usepackage{titlesec}
    \titleformat{\section}{\Large\bfseries}{\thesection}{1em}{}
    \titleformat{\subsection}{\large\bfseries}{\thesubsection}{1em}{}
    \titleformat{\subsubsection}{\normalsize\bfseries}{\thesubsubsection}{1em}{}
    \titlespacing*{\section}{0pt}{2.5ex plus 1ex minus .2ex}{1.5ex plus .2ex}
    \titlespacing*{\subsection}{0pt}{2ex plus 1ex minus .2ex}{1ex plus .2ex}
    \titlespacing*{\subsubsection}{0pt}{1.5ex plus 1ex minus .2ex}{0.8ex plus .2ex}

    % List formatting
    \usepackage{enumitem}
    \setlist[itemize]{topsep=0.5ex, itemsep=0.25ex, parsep=0pt}
    \setlist[enumerate]{topsep=0.5ex, itemsep=0.25ex, parsep=0pt}

    % Essential packages
    \usepackage{graphicx}
    \usepackage[table]{xcolor}
    \usepackage{booktabs}
    \usepackage{array}
    \usepackage{tabularx}
    \usepackage{longtable}
    \usepackage{amsmath}
    \usepackage{amssymb}
    \usepackage[normalem]{ulem}
    \rowcolors{2}{gray!10}{white}
    \renewcommand{\arraystretch}{1.3}

    % Cross-references
    \usepackage{hyperref}
    \usepackage{cleveref}

    % Custom commands
    \newcommand{\pilcrow}{\S}

    % Hyperref setup
    \hypersetup{
        colorlinks=true,
        linkcolor=blue,
        filecolor=magenta,
        urlcolor=cyan,
    }

    % Code listings
    \usepackage{listings}
    \definecolor{codebg}{RGB}{248,248,248}
    \definecolor{codeframe}{RGB}{200,200,200}
    \definecolor{codekeyword}{RGB}{0,0,180}
    \definecolor{codecomment}{RGB}{0,128,0}
    \definecolor{codestring}{RGB}{163,21,21}
    \lstset{
        basicstyle=\ttfamily\small,
        breaklines=true,
        breakatwhitespace=false,
        frame=single,
        framerule=0.5pt,
        rulecolor=\color{codeframe},
        backgroundcolor=\color{codebg},
        keepspaces=true,
        showstringspaces=false,
        tabsize=4,
        xleftmargin=1em,
        xrightmargin=1em,
        aboveskip=1em,
        belowskip=1em,
        keywordstyle=\color{codekeyword}\bfseries,
        commentstyle=\color{codecomment}\itshape,
        stringstyle=\color{codestring},
        literate={_}{{\textunderscore}}1
                 {-}{{\textminus}}1
                 {<}{{\textless}}1
                 {>}{{\textgreater}}1
                 {~}{{\textasciitilde}}1
                 {^}{{\textasciicircum}}1
    }

    % YAML language definition for listings
    \lstdefinelanguage{YAML}{
        keywords={true,false,null,y,n},
        keywordstyle=\color{blue}\bfseries,
        sensitive=false,
        comment=[l]{\#},
        morecomment=[s]{/*}{*/},
        commentstyle=\color{gray}\ttfamily,
        stringstyle=\color{red}\ttfamily,
        morestring=[b]'',
        morestring=[b]'',
    }

    % TOML language definition for listings
    \lstdefinelanguage{TOML}{
        keywords={true,false},
        keywordstyle=\color{blue}\bfseries,
        sensitive=true,
        comment=[l]{\#},
        commentstyle=\color{gray}\ttfamily,
        stringstyle=\color{red}\ttfamily,
        morestring=[b]'',
        morestring=[b]'',
    }

    % Dockerfile language definition for listings
    \lstdefinelanguage{Dockerfile}{
        keywords={FROM,RUN,CMD,LABEL,EXPOSE,ENV,ADD,COPY,ENTRYPOINT,VOLUME,USER,WORKDIR,ARG,ONBUILD,STOPSIGNAL,HEALTHCHECK,SHELL},
        keywordstyle=\color{blue}\bfseries,
        sensitive=false,
        comment=[l]{\#},
        commentstyle=\color{gray}\ttfamily,
        stringstyle=\color{red}\ttfamily,
        morestring=[b]'',
        morestring=[b]'',
    }

    % JSON language definition for listings
    \lstdefinelanguage{JSON}{
        keywords={true,false,null},
        keywordstyle=\color{blue}\bfseries,
        sensitive=true,
        commentstyle=\color{gray}\ttfamily,
        stringstyle=\color{red}\ttfamily,
        morestring=[b]'',
    }

    % Code listing float environment
    \usepackage{float}
    \newfloat{listing}{htbp}{lol}[section]
    \floatname{listing}{Listing}

    % Admonition boxes
    \usepackage{tcolorbox}
    \tcbuselibrary{skins,breakable}

    % Admonition style definitions
    \newtcolorbox{admonitionbox}[2][]{
        enhanced,
        breakable,
        colback=#2!5,
        colframe=#2!80!black,
        fonttitle=\bfseries,
        title=#1,
        left=2mm,
        right=2mm,
    }

    % Synthon tuple boxes
    \newtcolorbox{synthonbox}{
        enhanced,
        colback=blue!5,
        colframe=blue!75!black,
        fontupper=\large,
        center upper,
        boxrule=1pt,
        arc=4pt,
    }

    \begin{document}

    \title{Proof That 10 Is Solitary}
    \date{\today}

    \maketitle

    \subsection{Abstract}

    We approached this problem with an impression we would find a counterexample. Why not, after all? The small solitary numbers --- $1,2,...,9$ --- form a sparse set, and the abundancy index $9/5$ did not look especially forboding. Of course, in the end the proof spirals around a singular candidate and acquires a lock on its intended target. One by one, every alternative --- every single one --- is  nixed. And all this carnage from a constraint that seems too weak to land the fatal blows--- until you actually compute through each branch.

    \subsubsection{What the Problem Demands}

    Let $\sigma(n)$ denote the sum of all positive divisors of $n$. The \textbf{abundancy index} of $n$ is $\sigma(n)/n$. A positive integer $n$ is called \textbf{solitary} if it is the unique positive integer with its abundancy index. For $n = 10$:

    \[\sigma(10) = 1 + 2 + 5 + 10 = 18, \qquad \frac{\sigma(10)}{10} = \frac{9}{5}.\]

    The claim: \textbf{no integer $m \neq 10$ satisfies $\sigma(m)/m = 9/5$.}

    The function $\sigma$ is multiplicative: if $\gcd(a,b) = 1$, then $\sigma(ab) = \sigma(a)\sigma(b)$. For a prime power, $\sigma(p^e) = (p^{e+1} - 1)/(p - 1)$. If $m$ satisfies $\sigma(m)/m = 9/5$, then $5\sigma(m) = 9m$.

    All of this is standard. The difficulty is that $9/5$ is a rational number in lowest terms, and abundancy indices carry a rigid denominator-divisibility constraint. But the constraint alone is not enough --- it generates an infinite space of candidates when $m$ is factored as $2^a \cdot 5^b \cdot r$ with $\gcd(r,10) = 1$. We need to show that infinity collapses to finitely many candidates and that all but one die.

    \subsubsection{The Denominator Constraint and Why $b = 0$ Fails}

    The critical tool is the \textbf{denominator-divisibility lemma}: if $\sigma(r)/r = A/B$ with $\gcd(A,B) = 1$, then every prime divisor of $B$ divides $r$. Moreover, if $p \mid B$ and $p \mid r$, writing $r = p \cdot s$ gives

    \[\frac{\sigma(s)}{s} = \frac{A \cdot p}{B \cdot \sigma(p)},\]

    and the denominator of this new fraction (in lowest terms) must divide $s$.

    Now factor any candidate $m$ as $m = 2^a \cdot 5^b \cdot r$, where $a,b \ge 0$ and $\gcd(r,10) = 1$. Multiplicativity gives $\sigma(m) = \sigma(2^a)\,\sigma(5^b)\,\sigma(r)$.

    \textbf{Case $b = 0$.} Then $\sigma(5^0) = 1$, and $5\sigma(m) = 9m$ becomes

    \[5(2^{a+1} - 1)\,\sigma(r) = 9 \cdot 2^a \cdot r,\]

    yielding

    \[\frac{\sigma(r)}{r} = \frac{9 \cdot 2^a}{5(2^{a+1} - 1)}.\]

    The lower bound $\sigma(r)/r \ge 1$ forces $2^a \le 5$, so $a \in \{0,1,2\}$. The three candidates for $\sigma(r)/r$ in lowest terms are $9/5$, $6/5$, and $36/35$. Each denominator is divisible by $5$. But $r$ is coprime to $10$, so $r$ cannot contain the factor $5$ that the denominator-divisibility lemma demands. \textbf{No solution with $b = 0$.} This case closes cleanly.

    It is the next case where the proof earns its length.

    \subsubsection{The $b \ge 1$ Case: Abundance of Apparent Possibility}

    When $b \ge 1$, $\sigma(5^b) = (5^{b+1} - 1)/4$. The equation $5\sigma(m) = 9m$ expands to

    \[5 \cdot (2^{a+1} - 1) \cdot \frac{5^{b+1} - 1}{4} \cdot \sigma(r) = 9 \cdot 2^a \cdot 5^b \cdot r.\]

    Rearrange:

    \[\frac{\sigma(r)}{r} = \frac{36 \cdot 2^a \cdot 5^{b-1}}{(2^{a+1} - 1)(5^{b+1} - 1)}. \tag{1}\]

    Now impose $\sigma(r)/r \ge 1$. This single inequality does something remarkable: it collapses an infinite $(a,b)$ plane into a finite table. Enumerating all pairs $(a,b)$ with $a \ge 0$, $b \ge 1$ that satisfy (1):

    \begin{table}[htbp]
    \centering
    \small
    \rowcolors{1}{white}{white}
    \begin{tabularx}{\textwidth}{>{\centering\arraybackslash}X >{\centering\arraybackslash}X >{\centering\arraybackslash}X}
    \toprule
    \rowcolor{gray!25}\textbf{$a$} & \textbf{$b$} & \textbf{$\sigma(r)/r$ (lowest terms)} \\
    \midrule
    \rowcolors{1}{white}{gray!10}
    $0$ & $1$ & $3/2$ \\
    $0$ & $2$ & $45/31$ \\
    $0$ & $3$ & $75/52$ \\
    $0$ & $4$ & $1125/781$ \\
    $0$ & $5$ & $625/434$ \\
    $0$ & $6$ & $28125/19531$ \\
    $0$ & $7$ & $46875/32552$ \\
    $1$ & $1$ & $1/1$ \\
    \bottomrule
    \end{tabularx}
    \end{table}

    For $a = 0$ and $b \ge 8$, the right-hand side of (1) simplifies to

    \[\frac{\sigma(r)}{r} = \frac{36 \cdot 5^{b-1}}{5^{b+1} - 1} = \frac{36/5}{5 - 5^{-(b-1)}}.\]

    As $b \to \infty$, this approaches $36/25 = 1.44$ from above. Direct computation shows the value drops below $1$ for all $b \ge 8$. For $a \ge 1$ and $b \ge 2$, and for $a \ge 2$ and $b = 1$, the inequality $\sigma(r)/r \ge 1$ fails --- no further candidates.

    I want to pause on what just happened. We started with an equation that, in principle, could have admitted solutions for arbitrarily large $a$ and $b$ --- an infinite search. The inequality $\sigma(r)/r \ge 1$ killed all but eight candidate pairs. This is the first crossing point: the object pushes back against the naive expectation that the case analysis would be interminable. It is not. The abundancy index carries a hidden finiteness constraint that the inequality exposes.

    The eight that survive are the ones we must now kill.

    \subsubsection{Candidates with Immediate Denominator Failure}

    Four candidates die on first contact with the denominator-divisibility lemma. Since $\gcd(r,10) = 1$, any candidate whose $\sigma(r)/r$ has denominator divisible by $2$ or $5$ cannot correspond to any integer $r$.

    \begin{itemize}
        \item $(0,1)$: $3/2$, denominator $2$ --- dead.
        \item $(0,3)$: $75/52$, denominator $52 = 2^2 \cdot 13$ --- dead.
        \item $(0,5)$: $625/434$, denominator $434 = 2 \cdot 7 \cdot 31$ --- dead.
        \item $(0,7)$: $46875/32552$, denominator $32552 = 2^3 \cdot 13 \cdot 313$ --- dead.
    \end{itemize}

    The surviving candidates are $(0,2)$, $(0,4)$, $(0,6)$, and $(1,1)$. The last one is the solution $m = 10$: $\sigma(r)/r = 1$ forces $r = 1$, yielding $m = 2^1 \cdot 5^1 = 10$. The other three require deeper elimination.

    \subsubsection{The Hardest Cases}

    \subsubsection{Candidate $(0,2)$}

    Here $\sigma(r)/r = 45/31$ with prime denominator $31$. The lemma forces $31 \mid r$. Write $r = 31 \cdot s$. Then

    \[\frac{\sigma(s)}{s} = \frac{45 \cdot 31}{31 \cdot \sigma(31)} = \frac{45}{32}.\]

    The denominator $32 = 2^5$ must divide $s$, contradicting $\gcd(s,10) = 1$. Dead.

    The reasoning is swift, but notice why it works: the denominator-divisibility lemma is iterable. It does not merely reject a candidate --- it chases the contradiction into smaller integers. This candidate fails not because $45/31$ is impossible as an abundancy index, but because the denominator's prime factor $31$ triggers a cascade. The number $31$ is unremarkable; what matters is that $\sigma(31) = 32$ introduces a factor of $2$ into the denominator of the reduced fraction at the next level. The cascade is sensitive to the \emph{specific} prime, not just its size.

    \subsubsection{Candidate $(0,4)$}

    $\sigma(r)/r = 1125/781$. The denominator $781 = 11 \cdot 71$. We must consider both prime factors.

    \textbf{Branch 1:} $11 \mid r$. Write $r = 11 \cdot s$.

    \[\frac{\sigma(s)}{s} = \frac{1125 \cdot 11}{781 \cdot 12} = \frac{375}{284} = \frac{375}{2^2 \cdot 71}.\]

    Denominator now contains $2^2$. But $\gcd(s,10) = 1$ forbids $2 \mid s$ --- dead.

    \textbf{Branch 2:} $71 \mid r$. Write $r = 71 \cdot s$.

    \[\frac{\sigma(s)}{s} = \frac{1125 \cdot 71}{781 \cdot 72} = \frac{125}{88} = \frac{125}{2^3 \cdot 11}.\]

    Denominator $88 = 2^3 \cdot 11$ again forces a factor of $2$ into $s$ --- dead.

    Both branches die. The candidate is eliminated.

    There is a structural observation worth making explicit: the branching here is not parallel. The lemma does not tell us \emph{which} prime divisor of the denominator initiated the cascade --- we must try both. In every branch, the cascade terminates in a contradiction within one step. It is tempting to conjecture that this always happens for the solitary numbers, but I do not know whether it holds in general. The proof for $10$ works because the cascades are uniformly short.

    \subsubsection{Candidate $(0,6)$}

    $\sigma(r)/r = 28125/19531$. The denominator $19531$ is prime. Write $r = 19531 \cdot s$.

    \[\frac{\sigma(s)}{s} = \frac{28125 \cdot 19531}{19531 \cdot 19532} = \frac{28125}{19532} = \frac{28125}{2^2 \cdot 19 \cdot 257}.\]

    The denominator contains $2^2$ --- contradiction. Dead.

    \subsubsection{The Proof Closes}

    Every candidate has been examined:

    \begin{table}[htbp]
    \centering
    \small
    \rowcolors{1}{white}{white}
    \begin{tabularx}{\textwidth}{>{\centering\arraybackslash}X >{\centering\arraybackslash}X}
    \toprule
    \rowcolor{gray!25}\textbf{Candidate} & \textbf{Fate} \\
    \midrule
    \rowcolors{1}{white}{gray!10}
    $(0,1)$ & Immediate denominator failure ($2$) \\
    $(0,2)$ & Cascade contradiction via $\sigma(31)=32$ \\
    $(0,3)$ & Immediate denominator failure ($2$) \\
    $(0,4)$ & Both branches cascade to factor $2$ \\
    $(0,5)$ & Immediate denominator failure ($2$) \\
    $(0,6)$ & Cascade contradiction via $\sigma(19531)$ \\
    $(0,7)$ & Immediate denominator failure ($2$) \\
    $(1,1)$ & Survives: $m = 10$ \\
    \bottomrule
    \end{tabularx}
    \end{table}

    The only solution to $\sigma(m)/m = 9/5$ is $m = 10$. Therefore $10$ is solitary.

    The proof's engine is the iterated denominator-divisibility lemma, applied across a finite candidate table generated by the inequality $\sigma(r)/r \ge 1$. The inequality prunes infinity to finiteness; the lemma prunes finiteness to uniqueness. Both are necessary. Neither alone suffices.

    \subsubsection{What This Proof Leaves Open}

    The proof I expected to find would have been a direct application of a general theorem: "if $n$ satisfies condition $C$, then $n$ is solitary." No such theorem exists. The solitude of each small number is proved by a bespoke argument --- different denominators, different cascades, different branch structures. What unifies them is the denominator-divisibility lemma, but that lemma is a tool, not a theory.

    The proof for $10$ succeeds because the candidate table is small and every cascade terminates in a single step. For larger solitary numbers, the candidate tables grow and cascades can run longer. The solitude of $10$ is not deep in the way that, say, the solitude of a prime power is deep. It is deep in a different way: the proof works because $10 = 2 \cdot 5$ is exactly small enough that the finite set of candidates is manageable by hand. The structural question is whether there exists a single solitary number whose proof requires a cascade of length greater than $1$ --- a cascade where one iteration of the lemma does not suffice, and the contradiction is genuinely deferred.

    If such a number exists, its proof would be qualitatively different from the proof for $10$. The argument would need to track not just whether a denominator contains a forbidden prime but whether the cascade eventually forces one --- a property of the prime's orbit under $\sigma$, which is not well understood. The smallest candidate for such behavior might be a number whose abundancy index generates a denominator with a prime $p$ such that $\sigma(p)$ itself has a prime factor forbidden by coprimality --- and finding such a $p$ is not a matter of computation but of structural necessity.

    I do not know whether that number exists. The proof above does not need it to exist --- it only needs the eight candidates for $10$ to die. But the proof would be richer if the question were settled.

    \begin{center}
    \rule{0.5\textwidth}{0.4pt}
    \end{center}

    \emph{Structural type: $\langle D_{\infty};\ T_{\bowtie};\ R_{\leftrightarrow};\ P_{\pm};\ F_{\hbar};\ K_\text{slow};\ G_{\aleph};\ \Gamma_\text{seq};\ \⊙;\ H_2;\ n{:}m;\ \Omega_{\mathbb{Z}_2} \rangle$}


    \end{document}
